\section{Graph halving and lattice
polytopes}\label{graph-halving-and-lattice-polytopes}

Let \(\Gamma\) be a finite-rank residually connected incidence geometry,
of rank \(n\ge 3\). Designate types 0 and 1 as vertices and edges.
Assume their truncation is a simple graph Q, and that an edge is
incident with a higher object exactly when both its endpoints are
incident with that object. For a flag F of higher objects, let \(Q_F\)
be its common vertex-edge shadow. All \(Q_F\) are connected and contain
an edge.

UC.Geometry supplies the connected chamber graph of every residually
connected finite-rank geometry. Restricting a chamber gallery to any set
of at least two types shows that the corresponding truncation is
connected, because successive complete incidence subgraphs overlap.
Apply this inside the residue at F; chamber extension supplies an edge.
The edge-shadow hypothesis says \(Q_F\) is the induced graph on its
common vertices.

For a connected graph Q, its bipartite double cover D(Q) has vertices
\((v,\varepsilon )\), \(\varepsilon \in \{0,1\}\), with
\((v,\varepsilon )\) adjacent to \((w,1-\varepsilon )\) when v is
adjacent to w. It is connected when Q has an odd cycle, and has two
components otherwise: parity of walks proves both assertions. Choose a
component C of D(Q). Define \(L_C(\Gamma )\) as follows. Its first two
types are (v,0) and (v,1) in C. A higher object is a pair \((x,C_x)\),
where \(C_x\) is a component of \(D(Q_x)\) contained in C. Incidence of
the first two types is adjacency, incidence with a higher object is
membership in \(C_x\), and two higher objects are incident when their
old objects are incident and their chosen components meet.

The precise compatibility requirement is

\[
\text{if every }Q_x\ (x\in F)\text{ is nonbipartite, then }Q_F
\text{ is nonbipartite},
\tag{PC}
\]

for every nonempty higher flag F. It is automatic if the higher shadows
of every flag are nested, or if Q is bipartite.

\textbf{Theorem C.Graph.} The system \(L_C(\Gamma )\) is always a
geometry. It is residually connected exactly when (PC) holds. It is thin
exactly when \(\Gamma\) is thin. A chamber-transitive automorphism group
of \(\Gamma\) acts faithfully and chamber-transitively on
\(L_C(\Gamma )\) when Q is nonbipartite; when Q is bipartite its
index-two color-preserving subgroup does so.

\textbf{Proof.} For a flag \(\Phi =\{(x,C_x):x\in F\}\) of higher lifts,
write \(W=\cap C_x\). If any \(Q_x\) is bipartite, restriction of
\(C_x\) to \(D(Q_F)\) is one component. Any two such restrictions agree:
incidence makes them meet over the connected \(Q_{x,y}\), and hence
select the same parity class there and over \(Q_F\). Nonbipartite
constituents impose no restriction. Thus W is one component if any
constituent is bipartite, and all of \(D(Q_F)\) otherwise. In particular
W always contains an edge, and is connected exactly under (PC) for this
flag. For F empty use \(W=C\).

An edge of W and a chamber of \(\Gamma\) extending its base edge and F
give a chamber extending \(\Phi\). Every additional higher object has a
unique component containing the chosen lift of that edge. The same
extension works for a flag already containing one lifted vertex. If it
contains both, their base edge is unique by simplicity and incident with
all the old higher objects by the shadow condition. This proves the
geometry axiom.

The residue at (v,0) is the old residue \(\Gamma _v\): a neighbour (w,1)
corresponds to the unique edge \{v,w\}, and the component for an
incident higher object is uniquely specified by (v,0). Simplicity and
the shadow condition preserve all incidences. The other color is
identical. Hence residues of flags containing a lifted vertex are
connected. For a flag of higher objects only, its two-type truncation is
W; when W is connected, every other residual object attaches to it
through a chamber. This proves residual connectivity under (PC). If (PC)
fails, choose its higher flag; all constituents have unique lifts but W
has two components. Were the resulting geometry residually connected,
the truncation lemma inside its residue would make W connected. This
proves necessity without treating a truncation as a residue.

Panels containing both new vertex types correspond to old panels missing
the same higher type. A panel missing one new vertex type corresponds to
an old panel missing the edge type, at the retained endpoint. Old panels
missing the vertex type already have two objects by simplicity. These
correspondences prove both directions of thinness.

If Q is nonbipartite, an old chamber
\((v,\{v,w\},x_{2},\ldots ,x_{n-1})\) corresponds bijectively to the new
chamber with (v,0),(w,1) and the components containing them. Old
automorphisms act naturally on this bijection. If Q is bipartite, the
same bijection uses the old chambers whose chosen vertex is in one fixed
color class. A chamber-transitive group contains a color-exchanging
element, since it exchanges the endpoints of a base edge through a
chamber move. Its color-preserving subgroup therefore has index two and
is transitive on the selected chambers, with the same chamber bijection.
Faithfulness in both cases follows from the retained-object argument
below. \(\square\)

\subsection{The retained objects}\label{the-retained-objects}

For completeness, apply UC.Action to the precise object sets in this
construction. Let K be G when Q is nonbipartite and its color-preserving
subgroup otherwise. Take X to be the disjoint union of old vertices and
higher objects, and Y the object set of the halved geometry. The map \[
(v,\varepsilon)\mapsto v,\qquad (x,C_x)\mapsto x
\tag{C4}
\] is an equivariant surjection \(Y\to X\). For nonbipartite Q the
chosen C is all D(Q); for bipartite Q it contains exactly one lift of
every old vertex. In the latter case its intersection with the double
cover of each connected shadow \(Q_x\) is one of that double cover's two
components. Thus every higher object also has a lift.

Encode a new vertex by its fixed color tag and old vertex. Encode
\((x,C_x)\) by the old object x and the two subsets
\(\{v:(v,0)\in C_x\}\), \(\{v:(v,1)\in C_x\}\) of X, retaining its old
type as a tag. Components are determined by their vertex sets, so the
encoding is injective and equivariant. UC.Action gives \[
\ker(K\curvearrowright Y)=\ker(K\curvearrowright X)=1.
\tag{C5}
\] The last equality follows because an old edge is uniquely determined
by its two endpoints in the simple graph Q. Fixing X therefore fixes all
old objects. When \(\Gamma\) is thin and (PC) holds, the proved geometry
is thin and residually connected, and its K-action is
chamber-transitive. UC.Geometry then identifies K with the full
type-preserving group. This argument does not infer the geometry or its
flag properties merely from faithfulness.

Suppose now \(\Gamma\) is regular and thin, with vertex-changing marker
a, edge-changing marker b, and higher markers \(r_i\); assume a commutes
with the higher markers. Under (PC), Theorem C.Graph produces a regular
thin residually connected geometry. Directly on a base edge, its panel
involutions are \(b^a,b,r_i\): b changes one endpoint, aba changes the
other, and each \(r_i\) changes its higher object. UC.Geometry
identifies their group as the full type-preserving group H. Lemma
C.Gallery identifies the stabilizer of every base subflag with the
subgroup generated by its missing-type markers, giving IP and
identifying the actual coset incidence with \(L_C(\Gamma )\).

There is also an exact correspondence between the two parity conditions.
In the base residue with group \(U_J\) from (E1), the vertex stabilizer
is \(Q_J=\langle b,R_J\rangle\). Its vertex graph is bipartite precisely
when \(a\notin H_J\). If \(H_J\) is proper, its two cosets give a
coloring, since \(Q_J\le H_J\) and a crosses colors. Conversely the
action on the two colors supplies a character killing b and all \(r_j\)
and sending a to the nontrivial element, so it excludes a from \(H_J\).
In a globally nonbipartite geometry, higher flags reduce by transitivity
to base flags with shadows \(U_{E\setminus S}\). Condition (PC) says

\[
[E\setminus\{i\}\in\mathcal F\text{ for all }i\in S]
\Longrightarrow E\setminus S\in\mathcal F.
\tag{G1}
\]

An upward family containing E satisfies (G1) exactly when it is
principal: let T be the set of i whose complementary coatom belongs to
\(\mathcal F\); (G1) gives \(E\setminus T\in \mathcal F\), and every
member of \(\mathcal F\) contains \(E\setminus T\). Thus geometric parity
compatibility is the least-support condition of Theorem C.Endpoint. The
mixed overlap in that theorem has a different role: it says the two
adjacent vertices determine one edge. Indeed P and \(P^a\) stabilize the
endpoints; their intersection is the edge's pointwise stabilizer \(R_E\)
exactly when no other edge has those endpoints, using transitivity on
panels and the regular chamber identification.

\subsection{Halving lattice polytopes}\label{halving-lattice-polytopes}

\textbf{Corollary C.Halving.} Let P be a regular rank-n abstract
polytope, \(n\ge 3\), whose face poset is a lattice. Replacing one
endpoint marker a by aba gives a C-group H normal of index at most two
in Aut(P). Its maximal-parabolic coset system is a regular hypertope,
with full type-preserving group H. The diagram duplicates the edge
labels from b to the remaining markers, and the \(a'–b\) label is
m/gcd(m,2) for \(m=|ab|\) finite, or infinity for infinite m. The index
is two exactly when the corresponding vertex graph is bipartite; for a
final-end operation use the dual graph on facets.

\textbf{Proof.} In a lattice polytope an edge is the join of its two
distinct vertices. Thus two vertices determine at most one edge, and
that edge belongs to a higher face exactly when both vertices do. Higher
faces of a flag form a chain, so their shadows are nested and (PC) is
automatic. Theorem C.Graph and the panel-marker identification prove the
asserted geometry, full group, and IP without an
alternating-semiregular-polytope theorem. Equation (E2) and the coloring
argument give the index. Finally \(a'b=(ab)^{2}\) and a commutes with
all \(r_i\), so \(a'r_i\) is conjugate to \(br_i\); all other products
are unchanged. \(\square\)

